\documentclass[11pt]{article}
\usepackage[margin=0.8in]{geometry}
\usepackage{amsmath,amssymb,amsthm,microtype,hyperref}
\newtheorem{theorem}{Theorem}
\begin{document}
\begin{center}
{\Large Precision Checks for Proposed Collatz Residue Maps}\par\medskip
Collatz proof project\quad---\quad September 5, 2026
\end{center}
Let $T$ denote the shortcut Collatz map, and let
$S(n)=(3n+1)/2^{v_2(3n+1)}$ on positive odd integers.
We record elementary, kernel-checked tests relevant to reusing a recent
one-bit mixing reduction [1]. They concern specific residue statements,
not every argument in that work or the possibility of a corrected reduction.

\paragraph{Statement being checked.}
Section 5.2 of [1] presents a gap-step transition as a permutation modulo 32.
Theorem 4.2 uses compressed images of burst residues modulo $2^K$ and gives
a signed count difference $(-1)^K$. These statements require distinguishing
integer representatives from well-defined functions on residue classes.

\begin{theorem}[Lost precision in a gap step]
There is no function $f:\mathbb N\to\mathbb N$ such that
$T(n)\bmod32=f(n\bmod32)$ for every $n\equiv3\pmod4$.
Indeed, for every natural $q$,
\[
 T(3+32q)=5+48q.
\]
\end{theorem}
\begin{proof}
The input is odd, so $2T(3+32q)=3(3+32q)+1$.
In particular, 3 and 35 have the same input residue but their outputs
are 5 and 53, with residues 5 and 21 modulo 32.
A function of the input residue alone cannot give both values.
\end{proof}

\begin{theorem}[Compressed images need additional information]
The burst inputs 29 and 61 have the same residue modulo 32, but
\[
 S(29)=11\equiv3\pmod8,\qquad
 S(61)=23\equiv7\pmod8.
\]
\end{theorem}
\begin{proof}
Both inputs are 1 modulo 4. Their $3n+1$ values are
$88=8\cdot11$ and $184=8\cdot23$, with odd quotients as displayed.
\end{proof}

\paragraph{Canonical representatives do not repair orbit transitions.}
If the counts at depth five instead use the representatives
$0\leq n<32$, the gap-producing burst inputs are 9, 25, and 29.
Their compressed images are 7, 19, and 11. Hence
$C_3(5)=2$, $C_7(5)=1$, and $C_3(5)-C_7(5)=1$, whereas the printed
signed formula gives $(-1)^5=-1$. The absolute difference is still one;
this check does not refute that absolute-count statement.
Selecting representatives also does not justify applying their transition
table to arbitrary orbit values in the same classes.

\paragraph{Repair and scope.}
The gap formula shows exactly which information was discarded: the parity
of $q$ determines whether the output residue is 5 or 21. Retaining the
input modulo 64 resolves this example. Compressed steps require precision
appropriate to the number of divisions by two; a residue table needs an
explicit proof that its recorded inputs determine its outputs.
Neither these repairs nor balanced representative counts prove pointwise
orbit mixing. We therefore do not import the audited statements as a
coverage theorem for the current Lean proof program.

\paragraph{Verification.}
\begingroup\raggedright
Five declarations in \texttt{Collatz.ResidueMapAudit} prove the two numerical
witnesses, impossibility of the gap residue function, the quotient formula,
and the canonical depth-five counts. They use ordinary Lean proofs and
normalization, without a new mathematical axiom. Collatz remains unproved.\par\endgroup

\smallskip\begingroup\raggedright\noindent
[1] Edward Y. Chang, \emph{A structural reduction of the Collatz conjecture
to one-bit orbit mixing}, arXiv:2603.25753v1, Sections 4.2 and 5.2.
\href{https://arxiv.org/html/2603.25753v1}{Versioned primary source},
accessed September 5, 2026. No mathematical priority claim is made.\par\endgroup
\end{document}
